\bigskip\noindent{\Large\bfseries\color{apren} Test 3. Revisió global}\par\medskip
\iniciTest{espvecT3}{b,b,c,a,d,c,b,a,d,c}
\begin{enumerate}
\pregunta Sigui $\mathbb E$ l’espai de matrius reals $2\times2$ i $\mathbb F$ el conjunt de matrius simètriques de traça $0$. Aleshores $\mathbb F$ és:
\begin{enumerate}
\opcio{a}{Un subespai de dimensió $1$}
\opcio{b}{Un subespai de dimensió $2$}
\opcio{c}{Un subespai de dimensió $3$}
\opcio{d}{No és subespai}
\end{enumerate}\feedback
\pregunta Quin és el subespai generat per $(1,0)$ i $(0,1)$ a $\mathbb R^2$?
\begin{enumerate}
\opcio{a}{$\{0\}$}
\opcio{b}{$\mathbb R^2$}
\opcio{c}{La recta $y=x$}
\opcio{d}{Només els dos vectors donats}
\end{enumerate}\feedback
\pregunta Si $U+W$ és suma directa, què es compleix?
\begin{enumerate}
\opcio{a}{$U=W$}
\opcio{b}{$U+W=\{0\}$}
\opcio{c}{$U\cap W=\{0\}$}
\opcio{d}{$U\cap W=U$}
\end{enumerate}\feedback
\pregunta La fórmula de Grassmann afirma que:
\begin{enumerate}
\opcio{a}{$\dim(U+W)=\dim U+\dim W-\dim(U\cap W)$}
\opcio{b}{$\dim(U+W)=\dim U\cdot\dim W$}
\opcio{c}{$\dim(U\cap W)=\dim U+\dim W$}
\opcio{d}{$\dim V=0$}
\end{enumerate}\feedback
\pregunta Quin conjunt és base de $\mathbb R^2$?
\begin{enumerate}
\opcio{a}{$\{(1,0),(2,0)\}$}
\opcio{b}{$\{(0,0),(1,1)\}$}
\opcio{c}{$\{(1,1),(2,2)\}$}
\opcio{d}{$\{(1,0),(1,1)\}$}
\end{enumerate}\feedback
\pregunta Una recta vectorial de $\mathbb R^3$ té dimensió:
\begin{enumerate}
\opcio{a}{$0$}
\opcio{b}{$2$}
\opcio{c}{$1$}
\opcio{d}{$3$}
\end{enumerate}\feedback
\pregunta Un pla vectorial de $\mathbb R^3$ té dimensió:
\begin{enumerate}
\opcio{a}{$1$}
\opcio{b}{$2$}
\opcio{c}{$3$}
\opcio{d}{$4$}
\end{enumerate}\feedback
\pregunta Si $\dim V=n$, qualsevol base de $V$ té:
\begin{enumerate}
\opcio{a}{$n$ vectors}
\opcio{b}{$n+1$ vectors}
\opcio{c}{$2n$ vectors}
\opcio{d}{Un nombre arbitrari de vectors}
\end{enumerate}\feedback
\pregunta En general, $U+W$ és:
\begin{enumerate}
\opcio{a}{La unió de $U$ i $W$}
\opcio{b}{La intersecció de $U$ i $W$}
\opcio{c}{El producte cartesià de $U$ i $W$}
\opcio{d}{El conjunt de sumes $u+w$, amb $u\in U$ i $w\in W$}
\end{enumerate}\feedback
\pregunta Quin criteri detecta dependència lineal?
\begin{enumerate}
\opcio{a}{Que tots els vectors siguin no nuls}
\opcio{b}{Que la família tingui un sol vector no nul}
\opcio{c}{Que algun vector sigui combinació lineal dels altres}
\opcio{d}{Que els vectors pertanyin a un mateix cos}
\end{enumerate}\feedback
\end{enumerate}